\section{Fairlet Decomposition: a Top-down Approach on $\gamma$-HST}\label{sec:log-approx}

In this section, we provide a complete description of the first phase in our $(r,b)$-fair $k$-median algorithm (described in Section~\ref{sec:high-level}), namely our scalable ($r,b$)-fairlet decomposition algorithm. 

The first step in our algorithm is to embed the input point set into a $\gamma$-HST (for a value of $\gamma$ to be determined later in this section).
Once we build a $\gamma$-HST embedding $T$ of the input points $P\subseteq \mathbb{R}^d$, the question is how to partition the points into $(r,b)$-fairlets. 
We assume that each node $v\in T$ is augmented with extra information $N_r$ and $N_b$ respectively denoting the number of {\em red} and {\em blue} points in the subtree $T(v)$ rooted at $v$. 
%Moreover, for each collection of points $S$, we denote the set of red points and blue points respectively by $S_r$ and $S_b$.
%Furthermore, we assume that the update time for removing points from the data structure $T$ is $\tldO(1)$.

To compute the total cost of a fairlet decomposition, it is important to specify the clustering cost model (e.g., $k$-median, $k$-center). Here, we define $\cost_{\med}$ to denote the cost of a fairlet (or cluster) with respect to the cost function of $k$-median clustering: for any subset of points $S\subset \mathbb{R}^d$, 
$\cost_{\med}(S) := \min_{p\in S}\sum_{q\in S} d(p,q)$ where $d(p,q)$ denotes the distance of $p,q$ in $T$.

In this section, we design a fast fairlet decomposition algorithm with respect to $\cost_{\med}$.
\begin{theorem}\label{thm:median-main}
There exists an $\tldO(n)$ time algorithm that given an $O(r^5 + b^5)$-HST embedding $T$ of the point set $P$ and balance parameters $(r,b)$, computes an $O(r^3+b^3)$-approximate $(r,b)$-fairlet decomposition of $P$ with respect to $\cost_{\med}$ on $T$. 
\end{theorem}
Thus, by Theorem~\ref{thm:median-main} and the bound on the expected distortion of embeddings into HST metrics (Theorem~\ref{thm:tree-metric-distortion}), we can prove Theorem~\ref{thm:main-fairlet}.
\begin{proofof}{\bf Theorem~\ref{thm:main-fairlet}.}
We first embed the points into an $O(r^5 + b^5)$-HST $T$ and then perform the algorithm guaranteed in Theorem~\ref{thm:median-main}. By Theorem~\ref{thm:tree-metric-distortion}, the expected distortion of our embedding to $T$ is $O(d\cdot \gamma \cdot \log_{\gamma} n)$ and by Theorem~\ref{thm:median-main}, there exists an algorithm that computes an $O(r^3+b^3)$-approximate fairlet-decomposition of $P$ with respect to distances in $T$. Hence, the overall algorithm achieves $O(d\cdot (r^8+b^8)\cdot \log n)$-approximation. 

Since the embedding $T$ can be constructed in time $O(d\cdot n \cdot \log n)$ and the fairlet-decomposition algorithm of Theorem~\ref{thm:median-main} runs in near-linear time, the overall algorithm also runs in $\tldO(n)$.   
\end{proofof}


Before describing the algorithm promised in Theorem~\ref{thm:median-main}, we define a modified cost function $\cost_{\smed}$ which is a simplified variant of $\cost_{\med}$ and is particularly useful for computing the cost over trees. Consider a $\gamma$-HST embedding of $P$ denoted by $T$ and assume that $\script{X}$ is a fairlet decomposition of $P$. 
%For each fairlet $D\in \script{X}$, $\lca(D)$ denotes the {\em least common ancestor (lca)} of the points contained in $D$ in $T$.
Moreover, $h(D)$ denotes the height of $\lca(D)$ in $T$. For each fairlet $D$, $\cost_{\smed}(D)$ is defined as
\begin{align}\label{eq:fairlet-cost}
\cost_{\smed}(D) &:= \sum_{q\in D} d(\lca(D),q) = \Theta(\card{D} \cdot \gamma^{h(\lca(D))}).
\end{align}
In particular, $\cost_{\smed}(D)$ relaxes the task of finding ``the best center in fairlets'' and at the same time, its value is within a small factor of $\cost_{\med}(D)$.
\begin{claim}\label{clm:simple-median-cost}
Suppose that $T$ is a $\gamma$-HST embedding of the set of points $P$. For any $(r,b)$-fairlet $S$ of $P$, $\cost_{\med}(S) \leq \cost_{\smed}(S) \leq (r+b)\cdot \cost_{\med}(S)$ where the distance function is defined w.r.t. $T$.
\end{claim}

In the rest of this section we prove the following result which together with Claim~\ref{clm:simple-median-cost} imply Theorem~\ref{thm:median-main}.
\begin{lemma}\label{lem:simple-median-guarantee}
There exists a near-linear time algorithm that given an $O(r^5+b^5)$-HST embedding $T$ of the point set $P$ and balance parameters $(r,b)$, computes an $O(r^2+b^2)$-approximate $(r,b)$-fairlet decomposition of $P$ with respect to $\cost_{\smed}$ on $T$. 
\end{lemma}

%%%%%%%%%%%%% High level description %%%%%%%%%%%%%%%%%%%%%%%


\begin{lemma}\label{lem:constant-heavy-points}
For any tree $T$ with the root vertex $v$, the number of heavy points with respect to $v$ in the ($r,b$)-fairlet decomposition constructed by $\Call{\minHeavy}{v,r,b}$ is at most $O(r^2 + b^2)$ times the minimum number of heavy points in any valid ($r,b$)-fairlet decomposition of $T$. 
\end{lemma}

%%% Defer proof
%We prove Lemma~\ref{lem:constant-heavy-points} in Section~\ref{sec:heavy-points}. Here, we show that the overall algorithm, \Call{\fairletDec}\ constructs an $O(r^2 + b^2)$-approximate ($r,b$)-fairlet decomposition with respect to $\cost_{\smed}$.
%We defer the proof of Lemma~\ref{lem:simple-median-guarantee} to the longer version of the paper.

%%% Defer proof
\iffalse
\begin{proofof}{\bf Lemma~\ref{lem:simple-median-guarantee}.}
The proof is by induction on height of $v$ in $T$. The base case is when $v$ is a leaf node in $T$ and the algorithm trivially finds an optimal solution in this case. Suppose that the induction hypothesis holds for all vertices of $T$ at height $h-1$. Here, we show that the statement holds for the vertices of $T$ at height $h$ as well. 


Let $\opt$ denote an optimal ($r,b$)-fairlet decomposition of the points in $T(v)$ with respect to $\cost_{\smed}$. 
Next, we decompose $\opt$ into $\gamma^d + 1$ parts: $\set{\opt_i}_{i\in [\gamma^d]}$ and $\opt_{\heavy}$. 
For each $i\in[\gamma^d]$, $\opt_i$ denotes the set of fairlets in $\opt$ whose $\lca$ are in $T(v_i)$. Moreover, $\opt_{\heavy}$ denotes the set of heavy fairlets with respect to $v$ and $\heavy_{\opt}$ denotes the set of heavy points with respect to $v$ in $\opt$. Lastly, $\heavy^i_{\opt} := \heavy_\opt \cap T(v_i)$ denotes the set of heavy points with respect to $v$ in $\opt$ that are contained in $T(v_i)$. 
 
Let $\sol$ denote the solution returned by \Call{\fairletDec}{$v, r, b$}. Similarly, we decompose $\sol$ into $\gamma^d+1$ parts:  $\set{\sol_i}_{i\in [\gamma^d]}$ and $\sol_{\heavy}$. Moreover, $\heavy_{\sol}$ denotes the set of heavy points with respect to $v$ in $\sol$ and for each $i\in [\gamma^d]$, $\heavy^i_{\sol} := \heavy_\sol \cap T(v_i)$ denotes the set of heavy points with respect to $v$ in $\sol$ that are contained in $T(v_i)$.

%%% BOUND on heavy points in our solution
%By Lemma~\ref{lem:constant-heavy-points},
%\begin{align}\label{ineq:heavy-cost}
%\card{{\heavy}_\sol} \leq \eta_{\heavy} \cdot rb \cdot \card{{\heavy}_\opt},
%\end{align}
%where $\eta_\heavy$ is a constant. 
\begin{claim}\label{clm:aug-cost}
For each $i\in [\gamma^d]$, there exists an $(r,b)$-fairlet decomposition of $P_i\setminus \heavy_{\sol}^i$ of cost at most $\cost(\opt_i) + (|\heavy^i _\opt| + |\heavy^i_\sol| \cdot (r+b)) \cdot (r^2 + b^2) \cdot \gamma^{h-1}$ where $P_i$ is the set of points contained in $T(v_i)$.% and is $(r,b)$-balanced.
\end{claim}

Hence, by the induction hypothesis, for each $i\in [\gamma^d]$,
\begin{align}\label{ineq:children-cost}
\cost(\sol_i) 
&\leq  c\cdot (r^2 + b^2) \cdot (\cost(\opt_i) \nonumber \\ 
&+ (|\heavy^i _\opt| + |\heavy^i_\sol| \cdot (r+b)) \cdot (r^2 + b^2) \cdot \gamma^{h-1})
\end{align}
Next, we bound the cost of $\sol$ by Lemma~\ref{lem:constant-heavy-points} and~\eqref{ineq:children-cost} as follows: 
\begin{align*}
\cost(\sol) &= \cost(\sol_{\heavy}) + \sum_{i\in[\gamma^d]} \cost(\sol_{i}) \\
			 &\leq \eta_{\heavy}\cdot (r^2 + b^2) \cdot \cost(\opt_{\heavy}) \\
			 &\quad + c \cdot (r^2+b^2) \cdot ({\eta_{\heavy} (r+b)^5\over \gamma} \cdot \cost(\opt_{\heavy}) \\ 
			 &\quad + \sum_{i\in [k^d]}\cost(\opt_i))\\%\rhd\text{Claim~\ref{clm:aug-cost}}\\
			 &\leq c\cdot (r^2 + b^2) \cdot \cost(\opt) \\
			 &\quad + (\eta_\heavy - {c\over 2})\cdot (r^2+ b^2) \cdot \cost(\opt_{\heavy}) \\
			 &\qquad \rhd \text{By setting } \gamma := 2\eta_{\heavy}(r+b)^5\\
			 &\leq c\cdot (r^2 + b^2) \cdot \cost(\opt) \rhd c \geq 2\eta_\heavy
\end{align*}
\end{proofof}
\fi
%%%% End of defer proof


\subsection{Description of Step 1: Minimizing the Number of Heavy Points}\label{sec:heavy-points}
In this section, we show that \Call{\minHeavy} algorithm invoked by \Call{\fairletDec} finds an $O(r^2 + b^2)$-approximate solution of {\em Minimum Heavy Points} problem. %(as stated in Question~\ref{q:min-heavy-points}). %Note that throughout this section we assume that $r\geq b$.

The high-level overview of \Call{\minHeavy} is as follows. For any subset of points $D\subseteq P$, we can compute in $O(1)$ what the maximal size $(r,b)$-balanced subset of $D$ is: w.l.o.g. suppose that $N_r \geq N_b$ and $r\geq b$. If $N_r \leq {r\over b}\cdot N_b$, the collection is $(r,b)$-balanced. Otherwise, it suffices to greedily pick maximal size $(r,b)$-fairlets (see procedure \Call{\unbalancedPoints} for the formal algorithm). This simple observation implies a lower bound on the size of any optimal solution of \emph{Heavy Points Minimization} with respect to $v$ and we use this value to bound the approximation guarantee of \Call{\minHeavy} algorithm.
\begin{claim}\label{clm:minimum-point-removal}
\Call{\unbalancedPoints}{$N_r, N_b, r,b$} correctly computes the minimum number of points that is required to be removed from $N_r \cup N_b$ so that the remaining points become $(r,b)$-balanced. Moreover, the solution returned by the procedure only removes points form a single color class.
\end{claim}   
For each $i\in [\gamma^d]$, let $(\tilde{x}^i_r,\tilde{x}^i_b)$ be the output of \Call{\unbalancedPoints}{$N^i_r, N^i_b, r,b$}. 
\begin{corollary}\label{cor:trivial-lb}
Any $(r,b)$-fairlet decomposition of the points in $T(v)$ has at least $\sum_{i\in[\gamma^d]} \tilde{x}^i_r + \tilde{x}^i_b$ heavy points.
\end{corollary}

\begin{algorithm}[t]
\caption{\Call{\minHeavy}{$\set{N_r^i, N_b^i}_{i\in [\gamma^d]}, r, b, \gamma$}}
%: finds ($r,b$)-fairlets with minimum number of heavy points}\label{alg:min-heavy}
\begin{algorithmic}[1]
	\item[]
	\COMMENT{Stage $1$: lower bound on the number of heavy points} 
	\FORALL{non-empty children $i\in[\gamma^d]$ of $v$}
		\STATE $(x^i_r, x^i_b) \leftarrow \Call{\unbalancedPoints}{N^i_r, N^i_b, r, b}$ 
	\ENDFOR
	\STATE $y_r \leftarrow \sum_{i\in[\gamma^d]} x^i_r, y_b \leftarrow \sum_{i\in[\gamma^d]} x^i_b$
	\STATE $c_{\dom} \leftarrow \argmax_{c\in\set{r,b}} y_c, \bar{c_{\dom}} \leftarrow \set{r,b}\setminus c_{\dom}$
	
	\item[]
	\COMMENT{Stage $2$: add free points with color $\bar{c_\dom}$}
	\IF{$(\sum_{i\in[\gamma^d]} x^i_r, \sum_{i\in[\gamma^d]}x^i_b)$ is $(r,b)$-balanced}
		\STATE{break}
	\ENDIF
	\FORALL{non-empty children $i\in[\gamma^d]$ of $v$}
		\STATE $x^i_{\bar{c_{\dom}}} \leftarrow x^i_{\bar{c_\dom}} + \max(\Call{\extraPoints}{\bar{c_{\dom}}, N^i_r - x^i_r, N^i_b - x^i_b, r, b}, y_{c_{\dom}} - y_{\bar{c_{\dom}}})$	
	\ENDFOR

	\item[]
	\COMMENT{Stage $3$: add points of non-saturated $(r,b)$-fairlets}
	\IF{$(\sum_{i\in[\gamma^d]} x^i_r, \sum_{i\in[\gamma^d]}x^i_b)$ is $(r,b)$-balanced}
		\STATE{break} 
	\ENDIF

		\FORALL{non-empty children $i\in[\gamma^d]$ of $v$}
		\STATE $(n_r, n_b) \leftarrow \Call{\freeFairlet}{N^i_r - x^i_r, N^i_b - x^i_b, r, b}$ 
		\STATE $x^i_r \leftarrow x^i_r + n_r$, $x^i_b \leftarrow x^i_b + n_b$
	\ENDFOR	
	\RETURN $(\set{x^i_r, x^i_b}_{i\in[\gamma^d]})$
\end{algorithmic}
\end{algorithm}

\paragraph{Stage 1: minimum number of heavy points.}
If $\sum_{i\in[\gamma^d]} \tilde{x}^i_r$ red points together with $\sum_{i\in[\gamma^d]}\tilde{x}^i_b$ blue points form an $(r,b)$-balanced collection, then \Call{\minHeavy} technically terminates at the end of stage $1$ and the solution returned by \Call{\minHeavy} achieves the minimum possible number of heavy points. 
However, in general, the collection with $\sum_{i\in[\gamma^d]} \tilde{x}^i_r$ red points and $\sum_{i\in[\gamma^d]}\tilde{x}^i_b$ may not form an $(r,b)$-balanced collection. Next, we show that we can always pick at most $rb(\sum_{i\in[\gamma^d]} \tilde{x}^i_r + \tilde{x}^i_b)$ additional heavy points and keep both all subtrees rooted at children of $v$ and the set of heavy points $(r,b)$-balanced.  


\begin{algorithm}
\caption{\Call{\unbalancedPoints}{$N_r, N_b, r, b$}: returns the minimum number of points that are required to be removed so that $(N_r, N_b)$ become $(r,b)$-balanced.}\label{alg:min-unbalanced}
\begin{algorithmic}[1]
	%\Comment{assuming $r\geq b$}
	\IF{$N_r \geq N_b$} \RETURN {$(N_r - \floor{N_b \cdot {r\over b}}, 0)$}
	\ENDIF
	\RETURN $(0, N_b - \floor{N_r \cdot {r\over b}})$
\end{algorithmic}
\end{algorithm}


Another structure we will refer to in the rest of this section is {\em saturated ($r,b$)-fairlets}. A fairlet $D$ is a saturated $(r,b)$-fairlet if it has {\em exactly} $r+b$ points; $r$ points from color $c$ and $b$ points from color $\bar{c}$. 
\paragraph{Stage 2: Adding free points.}
If the ``must-have'' heavy points are not $(r,b)$-balanced, then one color is {\em dominant}. For a color class $c\in\set{r,b}$, a collection of points $S$ is $c$-dominant if $\card{S_c} \geq {r\over b}\cdot \card{S_{\bar{c}}}$. Moreover, the collection is {\em minimally-balanced} $c$-dominant if $S$ is $(r,b)$-balanced but it will be no longer $(r,b)$-balanced even if we remove a single point of color $\bar{c}$.

Let $c$ be the dominant color in the {heavy points}. Then, we inspect all children of $v$ and if there exits a child in which $\bar{c}$ is dominant, we borrow as many points of color $\bar{c}$ as we can (we need to keep the subtree $(r,b)$-balanced, see \Call{\extraPoints} procedure) till either the set of heavy points becomes $(r,b)$-balanced or all subtrees rooted at children of $v$ become minimally-balanced $c$-dominant. It is straightforward to show that at most ${b\over r}\cdot \card{S_c}$ points of color $\bar{c}$ will be borrowed from the children of $v$ in this phase.

\begin{algorithm}
\caption{\Call{\extraPoints}{$c, N_r, N_b, r, b$}: returns the maximum number of points of color $c$ that can be removed from the set $(N_r, N_b)$ such that they remain $(r,b)$-balanced.}\label{alg:extra-points}
\begin{algorithmic}[1]
%\Comment{assuming $r\geq b$}
	\IF{$N_c \leq N_{\bar{c}}$} \RETURN $0$ \ENDIF
	\RETURN $\ceil{N_{\bar{c}} \cdot {b\over r}}$
\end{algorithmic}
\end{algorithm}

\begin{lemma}\label{lem:structure-phase-2}
Suppose that the set of heavy points is $c$-dominant. If the set of heavy points is not $(r,b)$-balanced at the end of stage 2, then for each $i\in [\gamma^d]$, the set of points in the subtree rooted at $v_i$ is minimally-balanced $c$-dominant.
\end{lemma}
\begin{corollary}\label{cor:non-maximal}
Suppose that the set of heavy points is $c$-dominant. If the set of heavy points is not $(r,b)$-balanced at the end of stage 2, then for each $i\in [\gamma^d]$, the set of points in the subtree rooted at $v_i$ have an $(r,b)$-fairlet decomposition with at most one {non-saturated} $(r,b)$-fairlet. 
\end{corollary}

\paragraph{Stage 3: Non-saturated fairlets.}
Here, we show that we can increase the number of heavy points by at most a factor of $O(rb)$ %(assuming the whole data set is $(r,b)$-balanced) 
and make both the set of heavy points and the set of points in all subtrees rooted at children of $v$ $(r,b)$-balanced. Let $ns$ denote the total number of non-saturated fairlets in the subtree rooted at $v$. We consider two cases depending on the value of $ns$ and the total number of heavy points $N_\heavy$ that do not belong to any saturated fairlets (in particular, $\card{N_\heavy} \leq \sum_{i\in[\gamma^d]} \tilde{x}^i_r + \tilde{x}^i_b$): 

\paragraph{Case 1: $\boldsymbol{ns \leq b\cdot N_\heavy}$.}
If we add all non-saturated fairlets, since the rest of fairlets in the subtree rooted at children of $v$ are saturated $(r,b)$-balanced, then this ``extended'' collection of heavy points has to be $(r,b)$-balanced. Otherwise, the whole data set itself is not $(r,b)$-balanced which is a contradiction. Moreover, the total number of heavy points is at most $ns \times (r+b) = O(rb\cdot N_\heavy)$

\paragraph{Case 2: $\boldsymbol{ns > b\cdot N_\heavy}$.} 
Here we show that after adding at most $b \cdot N_{\heavy}$ non-saturated $(r,b)$-fairlets, the set of heavy points becomes $(r,b)$-balanced. Let $r_i$ and $b_i$ ($r_i\geq b_i$) specify the size of the non-saturated fairlet that belongs to the $i$-th child of $v$. Note that since all fairlets are $c$-dominant, $r_i$ denotes the number of points of color $c$.  

Moreover, in any  non-saturated $(r,b)$-fairlet, ${r_i \over b_i} < {r \over b}$, which implies that $r_i b \leq rb_i -1$. 
Let $Q$ denote the set of children of $v$ whose non-saturated fairlets are picked. After adding all points in these non-saturated fairlets, %of the subtree specified in $Q$ to the set of heavy points,
\begin{align}
\# \text{points of color } c 
&\leq N_{\heavy} + \sum_{j\in Q} {r_j} \leq N_{\heavy} + \sum_{j\in Q} ({r \over b}\cdot  b_j - {1\over b}) \leq {r\over b} \sum_{j\in Q} b_j \rhd\text{since $\card{Q} =b\cdot N_{\heavy}$},\label{eq:heavy-point-count-c}\\
\# \text{points of color } \bar{c} &\geq {r\over b} \sum_{j\in Q} b_j.\label{eq:heavy-point-count-bar-c}
\end{align}
Moreover, since in the beginning of the process the number of points of color $c$ is more than the number of points of color $\bar{c}$ and also in each non-saturated fairlest the number of points of color $c$ is more than the number of points of color $\bar{c}$, at the end of the process, in heavy points, the size of color $c$ is larger than the size of color $\bar{c}$.
Thus, by~\eqref{eq:heavy-point-count-c}~and~\eqref{eq:heavy-point-count-bar-c}, at the end of stage $3$, the extended heavy points has size $O(rb\cdot N_{\heavy})$ and is $(r,b)$-balanced as promised in Lemma~\ref{lem:constant-heavy-points}. %Hence, at the end of stage $3$, the set of heavy points picked by our algorithm is $(r,b)$-balanced an their size is not more than $O(r^2+b^2)$ times the size of the minimum possible number of heavy points as promised in Lemma~\ref{lem:constant-heavy-points}. 



\paragraph{Runtime analysis of \minHeavy.} Here we analyze the runtime of \Call{\minHeavy} which corresponds to step $1$ in \Call{\fairletDec}. Note that stage $1$ only requires $O(1)$ operations on the number of red and blue points in $T(v)$. Each of stage $2$ and stage $3$ requires $O(1)$ operations on the number of red and blue points in all non-empty children of $T(v)$. Although the number of children of $T(v)$ can be as large as $\gamma^d$, for each node $v$ in $T$, \Call{\minHeavy} performs $O(1)$ operations on the number of red and blue points in $T(v)$ exactly twice: when it is called on $v$ and the parent of $v$. Hence, in total $\Call{\minHeavy}$ performs $O(1)$ time on each node in $T$ which in total is $O(n)$.

%\begin{lemma}\label{lem:non-saturated}
%Suppose that $N_{\heavy}$ is the total number of heavy points at the end of phase 1. Then, at the end of Phase 3, the total number of heavy points is at most $b(r+b)N_{\heavy}$.
%\end{lemma} 
\begin{algorithm}
\caption{\Call{\freeFairlet}{$N_r, N_b, r, b$}: returns the non-saturated fairlet in a set with $(N_r, N_b)$ points.}\label{alg:non-saturated-fairlet}
\begin{algorithmic}[1]
%\Comment{assuming $r\geq b$}
	\IF{$N_r \leq N_{b}$} 
		\STATE $z_r \leftarrow N_r - \floor{N_r \over b}\cdot b$, $z_b \leftarrow N_b - \floor{N_r \over b}\cdot r$
	\ELSE{} 
		\STATE $z_b \leftarrow N_b - \floor{N_b \over b}\cdot b$, $z_r \leftarrow N_r - \floor{N_b \over b}\cdot r$
	\ENDIF
	\RETURN $(z_r, z_b)$
\end{algorithmic}
\end{algorithm}

